Reading technologies have significantly reshaped the way individuals consume and interact with textual content. While conventional electronic book (eBook) readers have democratized access to literature, they remain largely passive platforms, offering limited support for deeper cognitive engagement beyond basic features such as highlighting, bookmarking, and keyword search~\cite{b1}. In parallel, the advent of powerful large language models (LLMs), such as OpenAI’s ChatGPT~\cite{b2} and Meta’s LLaMA series~\cite{b3}, has introduced highly interactive paradigms that allow users to upload complete documents and receive artificial intelligence (AI)-generated summaries or responses to natural language queries.

These systems offer efficiency, yet they displace the cognitive effort paramount for critical reading. By giving answers rather than facilitating discovery, they risk reducing readers to passive recipients. LLM-based tools that process entire texts at once often reveal spoilers, disrupting narrative continuity, while traditional eReaders lack support for tracking complex threads. This dual shortfall—AI that does too much and eReaders that do too little—showcases the need for a system that leverages AI to augment, rather than supplant, the reading experience.

In response, this study introduces MeReader, an offline, privacy-preserving, and progress-aware eBook reader designed to resolve the limitations of current market alternatives.
While "progress gating" exists in commercial ecosystems like Amazon Kindle's "Ask This Book"~\cite{engadget_kindle,verge_kindle}, it is a closed feature with proprietary formats and cloud services, both of which are antithetical to the philosophy and the aim of MeReader.
There exists also, in the realm of PDFs, tools like the Adobe Acrobat AI Assistant~\cite{verge_pdfs} that tackle primarily document summarization but their focus is more on productivity rather than narrative eBook reading; treating a novel like a database and querying answers and/or summaries based on prompts without regards to automatic spoiler prevention.
Open source alternatives, such as the ever popular KOReader~\cite{github_koreader} with its AI Assistant plugin~\cite{github_koreader_plugin}, prioritize privacy, but rely on a "trailing window" of recent text (limited to $~100$k characters) which in turn leads to an amnesia, of sorts, in its retrieval and generation system; the inability to recall plot points from early chapters once the reader advances deep in the book.
MeReader amalgamates these capabilities by generalizing the spoiler-safety of Kindle into an open architecture (one that is auditable) while also maintaining the deep, full-book retrieval memory lacking in the KOReader plugin.

This study aims to advance digital reading technologies by mitigating cognitive constraints associated with electronic text consumption, particularly reduced spatial memory and diminished recall, through the development of a minimally intrusive, context-sensitive AI assistant. The proposed system aims to maintain narrative coherence and reader agency while improving comprehension through localized, progress-aware support mechanisms. 
The remainder of this paper is organized as follows. Section~\ref{sec:related_work} reviews related work on AI-augmented reading tools, semantic retrieval architectures, privacy-preserving inference methods, and progress-aware user modeling. Section~\ref{sec:system_design} presents the design and architecture of MeReader, including EPUB parsing, text embedding, and system integration. Section~\ref{sec:implement} describes the technical implementation and software stack, detailing model optimization, UI/UX considerations, and inference logic. Section~\ref{sec:evaluation} reports the results of empirical evaluation across representative reader tasks. Section~\ref{sec:conclusion} discusses the broader implications of the findings and system limitations and concludes with a summary of contributions, design novelties, and avenues for future development.